% Compile with: uplatex file.tex && dvipdfmx file.dvi
%           or: latexmk -uplatex file.tex
%\documentclass[uplatex,dvipdfmx]{jsarticle}
\documentclass[dvipdfmx,a4paper,11pt]{jsarticle}
\usepackage{amsmath, amssymb}
\usepackage{tikz}

\title{scenario\_two\_level\_perturbation}
\date{\today}

\begin{document}
\maketitle

2準位系の摂動論。$H=H_0+\lambda V$ で、$H_0$ の固有値が $E_a, E_b$、摂動 $V$ は非対角のみ。$\lambda$ は結合定数。

Ea, Eb, V, $\lambda$ をスカラーパラメータとして定義して

\[
  Ea, Eb, V, \lambda
\]

ハミルトニアンを $2\times2$ 行列として登録する。

$H = [Ea, \lambda V; \lambda V, Eb]$

\begin{equation}
  H = \begin{pmatrix} Ea & \lambda V \\ \lambda V & Eb \end{pmatrix}
  \label{eq:E1}
\end{equation}

$2\times2$ なので固有値は厳密に求まる。摂動級数はこの閉じた形を $\lambda$ で展開したものにほかならず、専用の「摂動展開」操作は要らない。

E1 の固有値を求めて

\begin{equation}
  \left\{ -\left(\frac{\sqrt{4 V^2 \lambda^2+Eb^2-2 Ea Eb+Ea^2}-Eb-Ea}{2}\right) , \frac{\sqrt{4 V^2 \lambda^2+Eb^2-2 Ea Eb+Ea^2}+Eb+Ea}{2} \right\}
  \label{eq:E2}
\end{equation}

下側の枝を式として登録する。

$(Ea + Eb - \sqrt{4 V^2 \lambda^2 + (Ea-Eb)^2})/2$

\begin{equation}
  (Ea + Eb - \sqrt{4 V^2 \lambda^2 + (Ea-Eb)^2})/2
  \label{eq:E3}
\end{equation}

まず1次まで展開する。摂動が非対角のみ（$\langle n|V|n\rangle=0$）なので、1次補正は消えるはず。

E3 を $\lambda$ について 0 のまわりで 1 次までテイラー展開して

\begin{equation}
  Ea
  \label{eq:E4}
\end{equation}

2次まで展開する。教科書の $E^{(2)}_n=|\langle m|V|n\rangle|^2/(E_n-E_m)$ が出るはず。

E3 を $\lambda$ について 0 のまわりで 2 次までテイラー展開して

\begin{equation}
  Ea-\frac{V^2 \lambda^2}{Eb-Ea}
  \label{eq:E5}
\end{equation}

上側の枝も同様に登録して展開する。符号が逆になり、2準位が互いに離れる（準位反発）。

$(Ea + Eb + \sqrt{4 V^2 \lambda^2 + (Ea-Eb)^2})/2$

\begin{equation}
  (Ea + Eb + \sqrt{4 V^2 \lambda^2 + (Ea-Eb)^2})/2
  \label{eq:E6}
\end{equation}

E6 を $\lambda$ について 0 のまわりで 2 次までテイラー展開して

\begin{equation}
  \frac{V^2 \lambda^2}{Eb-Ea}+Eb
  \label{eq:E7}
\end{equation}

検算。2つの準位の和はトレース $E_a+E_b$ に等しく、しかもこれは厳密に成り立つ。2次の寄与が打ち消し合えば級数が正しい。

$E5 + E7$

\begin{equation}
  E5 + E7
  \label{eq:E8}
\end{equation}

E8 を簡約して

\begin{equation}
  Eb+Ea
  \label{eq:E9}
\end{equation}

\end{document}